\documentclass[10pt, french]{beamer}
\usepackage[utf8]{inputenc}
\usepackage[T1]{fontenc}
\usepackage{babel}
\usepackage{amsmath, amssymb}
\usepackage{graphicx}
\usepackage{xcolor}

% Thème
\usetheme{Madrid}
\usecolortheme{default}

% Désactiver le pied de page par défaut (auteur/institution en bas à gauche)
\setbeamertemplate{footline}{}

% Informations
\title{Fight Fire with Fire}
\subtitle{Le regard des hacktivistes sur la désinformation sur les réseaux sociaux}
\author{Sébastien Léglise \& Allan Legrand}
\institute{ENSI Caen -- Avril 2026}
\date{Présentation (10 min)}

\begin{document}

% Diapositive 1 : Titre
\begin{frame}
\titlepage
\center{\tiny D'après Sharevski \& Kessell, SOUPS 2023}
\end{frame}

% Diapositive 2 : Plan
\begin{frame}{Plan de la présentation}
\tableofcontents
\end{frame}

\section{Contexte \& État de l'art}
\begin{frame}{Hacktivisme et éthique hacker}
\begin{itemize}
  \item \textbf{Hacktivisme} : détournement des outils techniques pour des causes militantes (Anonymous : OpKKK, OpFerguson, OpISIS).
  \item Arsenal : DDoS, \textit{doxing}, \textit{defacement}, mèmes, trolling.
  \item Vide post-2014 (arrestations) → comblé par des acteurs étatiques (\textit{troll farms}, \textit{sock puppets}).
  \item \textbf{Éthique hacker} (Levy, 1984) :
    \begin{enumerate}
      \item Toute information doit être libre.
      \item Méfiance envers l'autorité, promotion de la décentralisation.
    \end{enumerate}
  \item La désinformation contredit ces deux postulats.
\end{itemize}
\end{frame}

\section{Problématique \& Méthodologie}
\begin{frame}{Questions de recherche et méthode}
\textbf{3 questions de recherche :}
\begin{enumerate}
  \item Comment les hacktivistes conceptualisent-ils l'écosystème de la désinformation ?
  \item Quelles ripostes jugent-ils appropriées ?
  \item Comment perçoivent-ils l'évolution future ?
\end{enumerate}
\vspace{0.3cm}
\textbf{Méthode :}
\begin{itemize}
  \item Entretiens qualitatifs ouverts (Zoom, 2022-2023, $\approx$1h).
  \item $n = 22$ hacktivistes américains (actifs avant 2014).
  \item Analyse thématique (7 thèmes, accord inter-codeurs $\kappa = .82$).
\end{itemize}
\end{frame}

\section{Résultat}

\begin{frame}{Position des hacktivistes face à la désinformation}
\begin{itemize}
	\item \textbf{Menace pour la démocratie} : Aucun n'approuve cette récupération ; la désinformation est vue comme une menace pour la vision démocratique d'Internet.
  \item \textbf{Conflit éthique} : Elle s'oppose au principe d'information libre en créant un désordre qui catalyse la polarisation.
  \item \textbf{Méfiance de l'autorité} : Souvent orchestrée par une autorité (fantôme ou étatique), elle s'oppose à la décentralisation chère aux hackers.
  \item \textbf{Détournement des outils} : Les techniques historiques (trolling, mèmes) sont retournées pour manipuler les émotions.
  \item \textbf{Responsabilité des plateformes} : Les réseaux sociaux facilitent ce chaos en privilégiant la monétisation de l'engagement.
\end{itemize}
\end{frame}

\begin{frame}{Méthodes techniques préconisées}
\begin{itemize}
	\item \textbf{Neutralisation} : Interventions ciblées via le \textit{deplatforming} (fermeture de comptes), le \textit{leaking} et le \textit{doxing} des financeurs.
  \item \textbf{Saturation informationnelle} : « Combattre le feu par le feu » en inondant les canaux de désinformation avec des informations véridiques.
  \item \textbf{L'exemple \#OpJane} : Utilisation de "désinformation plausible" pour saturer et neutraliser les bases de données anti-avortement au Texas.
  \item \textbf{Pression économique} : Révélation d'identités (\textit{doxing}) pour forcer les annonceurs à se désolidariser des influenceurs néfastes.
\end{itemize}
\end{frame}

\begin{frame}{Recommandations pour la littératie informationnelle}
\begin{itemize}
	\item \textbf{Réforme éducative} : Intégration de véritables programmes d'apprentissage de la pensée critique dès l'école.
	\item \textbf{Responsabilité numérique} : L'enseignement de la sécurité opérationnelle (OpSec) devrait devenir une responsabilité sociétale.
	\item \textbf{Bots de vérité} : Développement de bots automatisés pour contrer les fermes de trolls et aider à sourcer les fait.
	\item \textbf{Neutralité du ton} : Abandon du ton moralisateur type « \textit{cancel culture} » au profit d'une attitude impartiale, les faits n'ayant pas de propriétés politiques.
\end{itemize}
\end{frame}

\section{Conclusion}
\begin{frame}{Conclusion}
\begin{enumerate}
  \item Les hacktivistes reconnaissent la désinformation comme une menace directe pour leur éthique fondatrice.
  \item Ils préconisent des ripostes concrètes (doxing, deplatforming, saturation) \textbf{et} des interventions éducatives.
  \item L'étude comble une lacune : ces acteurs qui ont inventé les outils aujourd'hui détournés n'avaient pas encore été consultés.
  \item Risque de fragmentation partisane du mouvement hacktiviste à surveiller.
\end{enumerate}
\vspace{0.5cm}
\centering
\small Sharevski \& Kessell, SOUPS 2023 \\
Rapport : Léglise \& Legrand, ENSI Caen, Avril 2026
\end{frame}

\end{document}
